% =====================================================================
%  Informe del Proyecto 3 - Grupo 3 - Desarrollo de Videojuegos (ESPE)
%  Unidad 2: Tres en Raya con IA en JavaScript (Web)
%  Compilar:  pdflatex InformeTickTackToe.tex  (dos veces, para el índice)
% =====================================================================
\documentclass[12pt,a4paper]{article}

\usepackage[utf8]{inputenc}
\usepackage[T1]{fontenc}
\usepackage[spanish,es-tabla,es-noshorthands]{babel}
\usepackage{lmodern}
\usepackage[margin=2.5cm]{geometry}
\usepackage{graphicx}
\usepackage{xcolor}
\usepackage{booktabs}
\usepackage{longtable}
\usepackage{array}
\usepackage{enumitem}
\usepackage{listings}
\usepackage{float}
\usepackage{caption}
\usepackage{subcaption}
\usepackage{fancyhdr}
\usepackage{amssymb}
\usepackage{tikz}
\usetikzlibrary{shapes.geometric,shapes.multipart,arrows.meta,positioning,calc,babel}
\usepackage[hidelinks]{hyperref}

\graphicspath{{Img/}}

% ---------- colores y utilidades ----------
\definecolor{espe}{RGB}{0,102,51}
\definecolor{xcol}{RGB}{8,145,178}
\definecolor{ocol}{RGB}{219,39,119}
\definecolor{codebg}{RGB}{245,245,250}

% Captura con borde fino; el argumento opcional recibe las opciones de la imagen.
\newcommand{\captura}[2][width=0.99\linewidth]{%
 \begingroup\setlength{\fboxsep}{0pt}\setlength{\fboxrule}{0.4pt}%
 \fcolorbox{gray!50}{white}{\includegraphics[#1]{#2}}\endgroup}

\captionsetup{font=small,labelfont=bf,labelsep=period}
\captionsetup[sub]{font=footnotesize,labelfont=bf}

% Salida real de consola, guardada en Img/consola-*.txt.
\definecolor{consolabg}{RGB}{24,26,38}
\definecolor{consolafg}{RGB}{226,232,240}
\lstdefinestyle{consola}{basicstyle=\ttfamily\scriptsize\color{consolafg},
  backgroundcolor=\color{consolabg}, frame=single, rulecolor=\color{consolabg},
  framesep=4pt, xleftmargin=4pt, xrightmargin=4pt, aboveskip=0pt, belowskip=0pt,
  breaklines=true, breakindent=1.2em, columns=fullflexible, keepspaces=true}

\lstdefinelanguage{JavaScript}{
  morekeywords={const,let,var,function,return,if,else,for,of,in,while,class,constructor,this,new,import,export,from,static,get,true,false,null,undefined,throw,async,await},
  sensitive=true, morecomment=[l]{//}, morecomment=[s]{/*}{*/}, morestring=[b]', morestring=[b]"}
\lstset{basicstyle=\ttfamily\footnotesize, backgroundcolor=\color{codebg},
  keywordstyle=\color{xcol}\bfseries, commentstyle=\color{gray}\itshape,
  stringstyle=\color{ocol}, frame=single, rulecolor=\color{gray!40},
  breaklines=true, columns=fullflexible, showstringspaces=false,
  literate={á}{{\'a}}1 {é}{{\'e}}1 {í}{{\'i}}1 {ó}{{\'o}}1 {ú}{{\'u}}1 {ñ}{{\~n}}1
    {¡}{{!`}}1 {✔}{{$\checkmark$}}1 {▶}{{$\blacktriangleright$}}1 {ℹ}{{\textit{i}}}1}

\pagestyle{fancy}
\fancyhf{}
\fancyhead[L]{\small Desarrollo de Videojuegos -- Unidad 2}
\fancyhead[R]{\small Grupo N.º 3}
\fancyfoot[C]{\thepage}
\renewcommand{\headrulewidth}{0.4pt}
\setlength{\headheight}{15pt}

\begin{document}
\sloppy

% =====================================================================
%  CARÁTULA
% =====================================================================
\begin{titlepage}
  \centering
  \includegraphics[width=0.55\textwidth]{logoINFESPE.png}\par\vspace{1cm}
  {\Large\bfseries UNIVERSIDAD DE LAS FUERZAS ARMADAS -- ESPE\par}
  \vspace{0.3cm}
  {\large DEPARTAMENTO DE CIENCIAS DE LA COMPUTACIÓN\par}
  \vspace{0.2cm}
  {\large CARRERA DE INGENIERÍA EN SOFTWARE\par}
  \vspace{1.2cm}
  {\large\bfseries ASIGNATURA: DESARROLLO DE VIDEOJUEGOS\par}
  \vspace{0.3cm}
  {\large UNIDAD 2\par}
  \vspace{1.2cm}
  {\color{espe}\rule{\linewidth}{1pt}}\par\vspace{0.3cm}
  {\LARGE\bfseries Tres en Raya con Inteligencia Artificial\par}
  \vspace{0.2cm}
  {\large Proyecto 3: Tic-Tac-Toe (MediumAI, UnbeatableAI y Minimax) -- JavaScript (Web)\par}
  \vspace{0.3cm}
  {\color{espe}\rule{\linewidth}{1pt}}\par
  \vspace{1.2cm}
  {\large DOCENTE: ING. CÉSAR VILLACÍS SILVA, MSC.\par}
  \vspace{0.8cm}
  {\large\bfseries ESTUDIANTES:\par}
  \vspace{0.2cm}
  {\large Diana Guerra\par Paul Moreno\par Germán Cáceres\par}
  \vspace{0.8cm}
  {\large\bfseries GRUPO N.º 3\par}
  \vfill
  {\large FECHA: NOVIEMBRE 2026\par}
\end{titlepage}

\tableofcontents
\newpage

% =====================================================================
\section{Presentación del proyecto}
% =====================================================================
Se desarrolló el videojuego \emph{Tres en Raya con Inteligencia Artificial} como
aplicación web en JavaScript. A diferencia de la versión de la Unidad~1, que
separaba un backend en Go de un frontend en Angular, en esta versión toda la
lógica del juego se ejecuta en el navegador como módulos ES, sin frameworks ni
dependencias externas. Un servidor estático mínimo en Node.js entrega los
archivos, porque los navegadores no cargan módulos ES desde \texttt{file://}.

\begin{table}[H]
\centering
\caption{Modalidades de inteligencia artificial implementadas}
\begin{tabular}{lll}
\toprule
\textbf{Modalidad} & \textbf{Tecnología} & \textbf{Estado}\\
\midrule
MediumAI & JavaScript (Web) & \textbf{Implementado}\\
UnbeatableAI & JavaScript (Web) & \textbf{Implementado}\\
UnbeatableAI con Minimax & JavaScript (Web) & \textbf{Implementado}\\
\bottomrule
\end{tabular}
\end{table}

Se migraron las tres variantes del proyecto C\#
\texttt{Tic-Tac-Toe-Game-Windows-Professor}, desarrollado originalmente con
Windows Forms. Se conservaron los pesos, el orden de evaluación de las reglas,
los mensajes de cada variante y, en Minimax, las puntuaciones y los criterios de
desempate. La única adaptación del muestreo aleatorio de UnbeatableAI se
describe en la sección~\ref{sec:hallazgos}. Las versiones de referencia son:

\begin{table}[H]
\centering
\caption{Versiones C\# utilizadas como referencia}
\small
\begin{tabular}{p{2.3cm}p{6.6cm}p{5cm}}
\toprule
\textbf{Variante} & \textbf{Carpeta} & \textbf{Clase}\\
\midrule
UnbeatableAI & \texttt{Tic-Tac-Toe-Versión-IA-Fuerte} & \texttt{UnbeatableAIController}\\
MediumAI & \texttt{Tic-Tac-Toe-Versión-IA-Debil} & \texttt{MediumAIController}\\
MinMax & \texttt{Tic-Tac-Toe-Version-MiniMax} & \texttt{UnbeatableAIController}\textsuperscript{*}\\
\bottomrule
\end{tabular}\\[2pt]
{\footnotesize \textsuperscript{*}Misma clase, pero implementada con el algoritmo Minimax en lugar del método numérico.}
\end{table}

% =====================================================================
\section{Especificación de Requerimientos (IEEE 830)}
% =====================================================================

\subsection{Introducción}

\subsubsection{Propósito}
Se especifican los requisitos funcionales y no funcionales del videojuego web
\emph{Tres en Raya con IA} en JavaScript, con base en IEEE 830-1998~\cite{ieee830}.
La especificación establece el comportamiento del sistema, sus interfaces y los
criterios de verificación de las funcionalidades implementadas.

\subsubsection{Alcance}
El sistema permite que un usuario juegue al Tres en Raya contra un
\emph{Jugador Controlado por la Computadora} (JCC) desde un navegador. La JCC
dispone de tres modalidades: \textbf{MediumAI} (1), \textbf{UnbeatableAI} (2) y
\textbf{UnbeatableAI con Minimax} (3). Las dos primeras usan pesos 0/1/3; Minimax
usa caracteres en una matriz 3$\times$3. El sistema comprende la gestión de la
partida, la alternancia de turnos, la detección del resultado y su presentación
animada en la misma página.

\subsubsection{Definiciones, acrónimos y abreviaturas}
\begin{description}[leftmargin=2.6cm,style=nextline]
  \item[JCC] Jugador Controlado por la Computadora (marca «O»).
  \item[Usuario] Jugador humano (marca «X»); siempre empieza.
  \item[mArrayMoves] Vector $V_k$ de 9 posiciones con los pesos 0/1/3.
  \item[Módulo ES] Archivo JavaScript que exporta e importa con \texttt{export} e \texttt{import}.
  \item[Evento] Registro de lo que el usuario ve: jugada de X, jugada de O o mensaje de resultado.
\end{description}

\subsubsection{Referencias}
\begin{itemize}
  \item IEEE Std 830-1998, \emph{Recommended Practice for Software Requirements Specifications}.
  \item Villacís, C. et al., \emph{Optimización del Juego del Tres en Raya con
        Niveles de Dificultad, Utilizando Técnicas Heurísticas de IA}, ESPE.
  \item Material de clase: \emph{Caso 1}, \emph{Caso 2} y \emph{Caso 3 TTT}.
  \item Código C\# de referencia: \texttt{Tic-Tac-Toe-Game-Windows-Professor} (IA fuerte, IA débil y MiniMax).
  \item Informe del Proyecto 3 de la Unidad~1 (versión Go y Angular).
\end{itemize}

\subsection{Descripción general}

\subsubsection{Perspectiva del producto}
La aplicación separa la lógica de la presentación dentro del mismo navegador.
La carpeta \texttt{src/game} contiene las reglas y las decisiones de IA migradas
desde C\#, sin acceso al DOM, de modo que también se ejecutan en Node.js para
las pruebas. La carpeta \texttt{src/ui} dibuja el tablero, recibe los clics y
reproduce los eventos que deja cada jugada.

\subsubsection{Funciones del producto}
Se permite seleccionar la modalidad de IA, iniciar o reiniciar una partida,
marcar una casilla libre, observar la respuesta de la JCC y consultar el
resultado, el marcador por modalidad, la bitácora de jugadas y el inspector de
pesos de las modalidades numéricas. MediumAI y UnbeatableAI emiten «¡Ganó X!»,
«¡Ganó O!» y «¡Empate!» con título «Resultado»; Minimax emite «You Win!»,
«AI Wins!» e «It's a Draw!» con título «Tic Tac Toe». El inspector numérico se
oculta en Minimax, porque ese algoritmo no usa pesos.

\subsubsection{Características de los usuarios}
Niños de 7 a 11 años (público objetivo del juego original) y estudiantes.
No se requiere conocimiento técnico: basta con hacer clic o usar las teclas 1--9.

\subsubsection{Restricciones}
\begin{itemize}
  \item La migración debe conservar el comportamiento funcional de las variantes C\#.
  \item JavaScript estándar (ECMAScript 2022) con módulos ES, sin frameworks ni dependencias de terceros.
  \item Node.js~20 o superior, solo para el servidor estático y las pruebas.
  \item Navegadores modernos (Chrome, Edge, Firefox) en computadora y celular.
\end{itemize}

\subsubsection{Suposiciones y dependencias}
La aplicación se sirve por HTTP; con \texttt{npm start} queda disponible en
\texttt{http://localhost:5173}. El marcador se guarda en \texttt{localStorage}
del navegador; si no está disponible, dura lo que dure la pestaña.

\subsection{Requerimientos específicos}

\subsubsection{Requerimientos funcionales}
\begin{longtable}{p{1.3cm}p{9.2cm}p{2.6cm}}
\caption{Requerimientos funcionales}\\
\toprule
\textbf{ID} & \textbf{Descripción} & \textbf{Prioridad}\\
\midrule\endfirsthead
\toprule
\textbf{ID} & \textbf{Descripción} & \textbf{Prioridad}\\
\midrule\endhead
RF01 & El sistema mostrará un tablero 3$\times$3. La lógica usa índices 0--8; el teclado de la interfaz admite 1--9. & Alta\\
RF02 & El usuario podrá seleccionar MediumAI (1), UnbeatableAI (2) o Minimax (3). & Alta\\
RF03 & El usuario siempre inicia con la marca «X». & Alta\\
RF04 & Una casilla libre recibe X; una ocupada se ignora sin provocar otra jugada de la IA. Los índices fuera de 0--8 se rechazan. & Alta\\
RF05 & En UnbeatableAI, en su primera jugada la JCC tomará el centro (4) si está libre; si no, una esquina aleatoria (0, 2, 6, 8). & Alta\\
RF06 & UnbeatableAI aplicará el método numérico: suma 6 $\Rightarrow$ gana; suma 2 $\Rightarrow$ bloquea; suma 5 en diagonal $\Rightarrow$ caso trivial; dos lados adyacentes del usuario $\Rightarrow$ esquina intermedia (1+3$\to$0, 1+5$\to$2, 3+7$\to$6, 5+7$\to$8); en otro caso, casilla aleatoria. & Alta\\
RF07 & En UnbeatableAI, si la JCC está en el centro y la regla de suma 5 aplica, jugará un lado aleatorio entre 1, 5 y 7 (comportamiento de la versión C\# descrito en la sección~\ref{sec:hallazgos}). & Alta\\
RF08 & Después de cada jugada se comprobarán filas, columnas y diagonales; el tablero lleno sin ganador produce empate, con los mensajes de cada variante. & Alta\\
RF09 & Al terminar la partida el tablero se deshabilitará (equivalente a \texttt{grbTicTacToe.Enabled = false}). & Alta\\
RF10 & El usuario podrá iniciar un juego nuevo en cualquier momento. & Media\\
RF11 & El sistema mostrará un marcador de partidas ganadas, empatadas y perdidas por modo. & Media\\
RF12 & El inspector mostrará pesos 0/1/3 y sumas por línea en las IA numéricas, y se ocultará en Minimax. & Baja\\
RF13 & MediumAI tomará primero el centro libre; en otro caso buscará suma 6, luego suma 2 y finalmente una casilla libre aleatoria. & Alta\\
RF14 & Minimax explorará con O maximizador y X minimizador; puntuará $10-d$, $-10+d$ o 0, restaurará las casillas y conservará el primer máximo estricto en orden 0--8. & Alta\\
RF15 & Tras un final no habrá más movimientos; un juego nuevo o un cambio de nivel reiniciará el tablero y el turno humano. & Alta\\
\bottomrule
\end{longtable}

\subsubsection{Requerimientos no funcionales}
\begin{longtable}{p{1.3cm}p{2.9cm}p{9.2cm}}
\caption{Requerimientos no funcionales}\\
\toprule
\textbf{ID} & \textbf{Tipo} & \textbf{Descripción}\\
\midrule\endfirsthead
\toprule
\textbf{ID} & \textbf{Tipo} & \textbf{Descripción}\\
\midrule\endhead
RNF01 & Fidelidad & Cada modo conservará las reglas de su variante C\#. Minimax mantendrá puntuaciones, orden y desempates, sin poda alfa-beta ni nuevas heurísticas.\\
RNF02 & Rendimiento & La decisión de la IA tardará menos de 50 ms por jugada, sin contar las pausas de animación.\\
RNF03 & Usabilidad & Interfaz animada, con sonido opcional, bitácora y soporte de teclado (1--9, N y S).\\
RNF04 & Accesibilidad & Respeta \texttt{prefers-reduced-motion} y usa etiquetas ARIA en el tablero y el resultado.\\
RNF05 & Portabilidad & Funciona en navegadores modernos y se adapta a celular, tableta y escritorio sin desplazamiento horizontal.\\
RNF06 & Mantenibilidad & Cada modo vive en su propio módulo y expone los mismos métodos (\texttt{playUser}, \texttt{playNPC}, \texttt{weights}, \ldots).\\
RNF07 & Fiabilidad & UnbeatableAI comprobará la disponibilidad antes del muestreo aleatorio restringido, para no congelar la pestaña.\\
RNF08 & Simplicidad & Cero dependencias de terceros en la aplicación, el servidor y las pruebas.\\
RNF09 & Calidad & Las pruebas de las tres modalidades deberán aprobar con \texttt{npm test}.\\
\bottomrule
\end{longtable}

% =====================================================================
\section{Modelo de Inteligencia Artificial}
% =====================================================================
En las dos IA numéricas, el tablero se representa como el vector $V_k$, con
$k = 0,\dots,8$, en el arreglo \texttt{arrayMoves} (\texttt{mArrayMoves} en C\#).
Cada casilla vale 0 (vacía), 1 (usuario) o 3 (JCC). Para cada una de las 8 líneas
se calcula la suma de sus tres casillas:

\begin{table}[H]
\centering
\caption{Interpretación de la suma por línea}
\begin{tabular}{cll}
\toprule
\textbf{Suma} & \textbf{Composición} & \textbf{Acción de la JCC}\\
\midrule
6 & 3 + 3 + 0 & Completar la línea y \textbf{ganar} (Caso 2)\\
2 & 1 + 1 + 0 & \textbf{Bloquear} al usuario (Caso 1)\\
5 (diagonal) & 1 + 3 + 1 & Caso trivial: esquina libre o lado aleatorio (Caso 3)\\
\bottomrule
\end{tabular}
\end{table}

Los pesos 1 y 3 hacen que cada suma relevante tenga una sola composición: con
1 y 2, por ejemplo, la suma 2 podría ser «X + X + vacía» u «O + vacía + vacía».
El C\# escribe 16 ramas \texttt{if/else if} para las sumas 6 y 2; en JavaScript
se expresan como dos recorridos de \texttt{WINNING\_LINES}, que tienen el mismo
orden (filas, columnas y diagonales) y devuelven la misma casilla:

\begin{lstlisting}[language=JavaScript]
// Suma 6 (3 + 3 + 0): la JCC completa la línea y gana (Caso 2).
for (const [a, b, d] of WINNING_LINES) {
  if (m[a] + m[b] + m[d] === 6) return this.pick(a, b, d);
}
// Suma 2 (1 + 1 + 0): la JCC bloquea al usuario (Caso 1).
for (const [a, b, d] of WINNING_LINES) {
  if (m[a] + m[b] + m[d] === 2) return this.pick(a, b, d);
}
\end{lstlisting}

\texttt{pick(a, b, d)} devuelve la primera casilla vacía de la línea y guarda su
suma en \texttt{heuristicSum}, como \texttt{mHeuristicSum} del original. Después
siguen, en el mismo orden que el C\#, la regla de suma 5 en cada diagonal, las
cuatro reglas de esquina por lados adyacentes y la jugada aleatoria.

\subsection{Variante MediumAI}
Se implementó \texttt{MediumAIController} mediante la migración de la variante
C\# de IA débil, con los pesos originales. El turno humano siempre inicia con X.
Después de cada jugada válida, el controlador comprueba las ocho líneas ganadoras
y el empate; al terminar, registra el mensaje del resultado y deshabilita el tablero.

La JCC (O) aplica las reglas en este orden:
\begin{enumerate}
  \item Si el centro (\texttt{CENTER\_INDEX = 4}) está vacío, lo elige.
  \item Si el centro está ocupado, busca una suma de 6 para completar una línea y ganar.
  \item Si no la hay, busca una suma de 2 para bloquear la primera amenaza de X.
  \item Si ninguna regla aplica, elige al azar una casilla vacía.
\end{enumerate}
Esta variante no incluye las reglas de suma 5 ni la estrategia de esquinas y lados
de UnbeatableAI.

\subsection{Variante UnbeatableAI con Minimax}
Minimax selecciona la jugada con el mejor resultado posible frente a un adversario
que también busca su mejor resultado. Se migraron los métodos de
\texttt{UnbeatableAIController.cs} y el flujo de \texttt{frmTicTacToe.cs} de la
variante \texttt{Tic-Tac-Toe-Version-MiniMax}~\cite{minimax-docente}.

\texttt{getBestMove} recorre 0--8 y calcula fila $\lfloor i/3\rfloor$ y columna
$i\bmod3$. Simula O en cada espacio libre, llama a \texttt{minimax(board, 0, false)}
y restaura el espacio. Actualiza la mejor jugada solo cuando
\texttt{score > bestScore}: conserva la primera de igual puntuación. Sin casillas
libres devuelve $-1$. \texttt{minimax} devuelve $10-d$ si gana O, $-10+d$ si gana
X y 0 en empate; en otro caso, O maximiza y X minimiza. Los límites iniciales son
los de \texttt{int.MinValue} e \texttt{int.MaxValue} de C\#.

\begin{table}[H]
\centering\small
\caption{Correspondencia de la variante Minimax C\# con JavaScript}
\begin{tabular}{p{6.3cm}p{6.6cm}}
\toprule Método o estructura C\# & Adaptación JavaScript\\\midrule
\texttt{char[3,3] board} & \texttt{board}: arreglo 3$\times$3 de \texttt{'X'}, \texttt{'O'} y \texttt{' '}\\
\texttt{GetBestMove / Minimax} & \texttt{getBestMove / minimax}\\
\texttt{EvaluateBoard / IsMovesLeft} & \texttt{evaluateBoard / isMovesLeft}\\
\texttt{ResetBoard} & \texttt{initializeData / initializeControls}\\
\texttt{OnButtonClick / AIMove} & \texttt{playUser / playNPC}, coordinados por \texttt{Form.click}\\
\texttt{MakeMove / CheckGameEnd} & \texttt{makeMove / checkGameEnd}\\
\texttt{CheckWin / IsBoardFull} & \texttt{checkWin / isBoardFull}\\
\texttt{MessageBox.Show} & \texttt{Board.showMessage}: evento con texto y título originales\\
\texttt{int.MinValue / int.MaxValue} & \texttt{-2147483648 / 2147483647}\\
\bottomrule
\end{tabular}
\end{table}

\section{Diagramas}
% =====================================================================
Las figuras~\ref{fig:uml-casos}, \ref{fig:uml-secuencia} y \ref{fig:uml-clases}
representan los casos de uso, la comunicación entre módulos y la estructura del
sistema implementado.

\subsection{Diagrama de casos de uso}
\begin{figure}[H]
\centering
\begin{tikzpicture}[
 uc/.style={ellipse,draw,fill=xcol!8,minimum width=4.4cm,minimum height=0.9cm,font=\small,align=center},
 sysuc/.style={uc,fill=ocol!8,minimum width=3.8cm},
 rel/.style={dashed,->},
 lbl/.style={font=\scriptsize,fill=white,inner sep=1pt},>=Stealth]
 \node[draw,rounded corners,minimum width=11.5cm,minimum height=10cm] (sys) at (1.15,0) {};
 \node[font=\bfseries\small,anchor=south] at (sys.north) {Tres en Raya IA (JavaScript)};
 \node[uc] (mode) at (-2,3.7) {Seleccionar modo:\\Medium / Unbeatable /\\Minimax};
 \node[uc] (new) at (-2,1.9) {Iniciar o reiniciar partida};
 \node[uc] (move) at (-2,0) {Jugar turno humano};
 \node[uc] (score) at (-2,-1.9) {Consultar marcador};
 \node[uc] (inspect) at (-2,-3.7) {Ver inspector de IA};
 \node[sysuc] (check) at (4.7,2.6) {Comprobar final};
 \node[sysuc] (ai) at (4.7,0) {Responder con IA};
 \node[sysuc] (result) at (4.7,-2.6) {Presentar resultado};
 \begin{scope}[shift={(-6.7,0)},line width=0.9pt,xcol!70!black]
  \draw[fill=xcol!8] (0,0.8) circle (0.27);
  \draw (0,0.53) -- (0,-0.3) (-0.42,0.2) -- (0.42,0.2) (0,-0.3) -- (-0.33,-0.9) (0,-0.3) -- (0.33,-0.9);
 \end{scope}
 \node[font=\small] at (-6.7,-1.2) {Usuario};
 \coordinate (usr) at (-6.1,0);
 \foreach \n in {mode,new,move,score,inspect} \draw (usr) -- (\n.west);
 \draw[rel] (move.north east) -- node[lbl]{$\ll$include$\gg$} (check.south west);
 \draw[rel] (ai.west) -- node[lbl]{$\ll$extend$\gg$} (move.east);
 \draw[rel] (result.north west) -- node[lbl]{$\ll$extend$\gg$} (move.south east);
\end{tikzpicture}
\caption{Casos de uso de las tres modalidades. La respuesta de la IA extiende el
turno si la jugada fue aceptada y la partida continúa; el resultado aparece al
terminar. La lógica forma parte del sistema, no es un actor externo.}
\label{fig:uml-casos}
\end{figure}

\subsection{Diagrama de secuencia: marcar casilla con Minimax}
\begin{figure}[H]
\centering\resizebox{\textwidth}{!}{%
\begin{tikzpicture}[font=\small,>=Stealth,
 obj/.style={draw,fill=xcol!10,minimum width=2.6cm,minimum height=1.2cm,align=center},
 msg/.style={align=center,fill=white,inner sep=1pt},
 self/.style={right,align=left}]
 \node[obj] (a) at (0,0) {:Usuario};
 \node[obj] (b) at (3.8,0) {:app.js\\(interfaz)};
 \node[obj] (c) at (7.6,0) {:Form};
 \node[obj] (d) at (11.4,0) {:Minimax-\\AIController};
 \node[obj] (e) at (15.2,0) {:Board};
 \foreach \n in {a,b,c,d,e} \draw[dashed,gray] (\n.south) -- ++(0,-14.1);
 \foreach \x/\ini/\fin in {3.8/-1.5/-13.9,7.6/-2.4/-10.6,11.4/-3.3/-5.9,11.4/-7.4/-10.4}
  \draw[fill=xcol!12,draw=xcol!60!black] (\x-0.09,\ini) rectangle (\x+0.09,\fin);
 \draw[gray,rounded corners=1pt] (7.0,-6.3) rectangle (17.6,-10.8);
 \node[anchor=north west,draw=gray,fill=white,font=\scriptsize,inner sep=2.5pt] at (7.0,-6.3) {\textbf{opt} [jugada aceptada y sin final]};
 \draw[->] (0,-1.5) -- node[msg,above=1pt]{clic o tecla $i$} (3.8,-1.5);
 \draw[->] (3.8,-2.4) -- node[msg,above=1pt]{click($i$)} (7.6,-2.4);
 \draw[->] (7.6,-3.3) -- node[msg,above=1pt]{playUser(board,$i$)} (11.4,-3.3);
 \draw[->] (11.4,-4.0) -- node[msg,above=1pt]{setText($i$, X)} (15.2,-4.0);
 \draw[->] (11.4,-4.5) -- ++(0.8,0) -- node[self]{checkGameEnd} ++(0,-0.6) -- (11.4,-5.1);
 \draw[dashed,->] (11.4,-5.8) -- node[msg,above=1pt]{true} (7.6,-5.8);
 \draw[->] (7.6,-7.4) -- node[msg,above=1pt]{playNPC(board)} (11.4,-7.4);
 \draw[->] (11.4,-7.8) -- ++(0.8,0) -- node[self]{getBestMove,\\minimax} ++(0,-0.8) -- (11.4,-8.6);
 \draw[->] (11.4,-9.2) -- node[msg,above=1pt]{setText($j$, O)} (15.2,-9.2);
 \draw[->] (11.4,-9.6) -- ++(0.8,0) -- node[self]{checkGameEnd} ++(0,-0.6) -- (11.4,-10.2);
 \draw[->] (3.8,-11.4) -- node[msg,above=1pt,pos=0.75]{takeEvents()} (15.2,-11.4);
 \draw[dashed,->] (15.2,-12.1) -- node[msg,above=1pt,pos=0.25]{eventos: X, O y MessageBox} (3.8,-12.1);
 \draw[->] (3.8,-12.5) -- ++(0.8,0) -- node[self]{replay: animar X,\\«pensando», O y resultado} ++(0,-0.8) -- (3.8,-13.3);
 \draw[dashed,->] (3.8,-13.8) -- node[msg,above=1pt]{tablero y tarjeta} (0,-13.8);
\end{tikzpicture}}
\caption{Jugada del nivel 3. \texttt{Form.click} solo despacha \texttt{playNPC}
si la jugada humana fue aceptada, y \texttt{playNPC} no actúa si la partida ya
terminó. En los niveles 1 y 2 cambia el controlador, pero el flujo es el mismo.}
\label{fig:uml-secuencia}
\end{figure}

\subsection{Diagrama de clases}
\begin{figure}[H]
\centering\resizebox{\textwidth}{!}{%
\begin{tikzpicture}[font=\scriptsize,>=Stealth,
 cls/.style={draw,rectangle split,rectangle split parts=3,rectangle split part fill={xcol!14,white,white},text width=5.6cm,align=left,inner sep=3pt,anchor=north},
 real/.style={dashed,-{Triangle[open,length=2.6mm,width=2.2mm]}}]
 \node[cls] (form) at (0,8) {\makebox[\linewidth]{\textbf{Form}}
  \nodepart{second}+level: number\\+controlPlayer: boolean\\+board: Board\\+medium: MediumAIController\\+unbeatable: UnbeatableAIController\\+minimax: MinimaxAIController
  \nodepart{third}+get controller(): controlador\\+setLevel(level)\\+initializeGame()\\+click(index)};
 \node[cls,rectangle split part fill={gray!15,white,white},text width=4.1cm] (app) at ([xshift=-7.4cm]form.north) {\makebox[\linewidth]{$\ll$módulo$\gg$ \textbf{ui/app.js}}
  \nodepart{second}state: form, view, busy, scores
  \nodepart{third}clickCell(i)\\replay(events)\\renderAll()\\selectMode(level)};
 \node[cls,text width=4.4cm] (board) at (7.3,8) {\makebox[\linewidth]{\textbf{Board}}
  \nodepart{second}+buttons: string[9]\\+enabled: boolean\\+events: Event[]
  \nodepart{third}+setText(index, mark, who)\\+clearTexts()\\+showMessage(text, caption)\\+takeEvents(): Event[]};
 \node[cls,rectangle split part fill={ocol!12,white,white}] (ctrl) at (0,3.2) {\makebox[\linewidth]{$\ll$interfaz implícita$\gg$ \textbf{Controlador}}
  \nodepart{second}Sin atributos
  \nodepart{third}+initializeData()\\+initializeControls(board)\\+playUser(board, index): boolean\\+playNPC(board)\\+get flag(): string\\+weights(): number[9]};
 \node[cls] (strong) at (-7.1,-1.3) {\makebox[\linewidth]{\textbf{UnbeatableAIController}}
  \nodepart{second}+currentPlayerFlag: Player\\+moveCount: number\\+currentGameState: GameState\\+boardState: Player[9]\\+arrayMoves: number[9]\\+heuristicSum: number
  \nodepart{third}+initializeGameData()\\+checkWinner(board): GameState\\+heuristicNumericalMethod(): number\\+pick(a, b, d): number\\+randomFreeAmong(...allowed): number\\+generateRandomNPCMove(): number};
 \node[cls] (weak) at (0,-1.3) {\makebox[\linewidth]{\textbf{MediumAIController}}
  \nodepart{second}+currentPlayerFlag: Player\\+moveCount: number\\+currentGameState: GameState\\+boardState: Player[9]\\+arrayMoves: number[9]
  \nodepart{third}+initializeGameData()\\+checkWinner(board): GameState\\+heuristicNumericalMethod(): number\\+firstEmpty(a, b, d): number\\+generateRandomNPCMove(): number};
 \node[cls] (minmax) at (7.1,-1.3) {\makebox[\linewidth]{\textbf{MinimaxAIController}}
  \nodepart{second}+board: string[3][3]\\+isHumanTurn: boolean\\+isGameOver: boolean
  \nodepart{third}+getBestMove(board): number\\+minimax(board, depth, isMaximizing): number\\+evaluateBoard(board): number\\+isMovesLeft(board): boolean\\+makeMove(board, index, player)\\+checkGameEnd(board): boolean\\+checkWin(player): boolean\\+isBoardFull(): boolean};
 \draw[dashed,->] (app.east |- form.center) -- node[above,align=center]{usa} (form.west |- form.center);
 \draw[{Diamond[fill=black]}-] (board.west -| form.east) -- node[above]{1} (board.west);
 \draw[dashed,->] (form.south) -- node[right]{controller} (ctrl.north);
 \draw[real] (strong.north) -- (ctrl.west);
 \draw[real] (weak.north) -- (ctrl.south);
 \draw[real] (minmax.north) -- (ctrl.east);
\end{tikzpicture}}
\caption{Estructura en UML. JavaScript no tiene interfaces: los tres controladores
exponen los mismos métodos y \texttt{Form} los usa sin distinguirlos. Los métodos
comunes del controlador se omiten en cada clase concreta.}
\label{fig:uml-clases}
\end{figure}

% =====================================================================
\section{Metodología SCRUM + XP}
% =====================================================================
El desarrollo se organizó en tres sprints, con Scrum como marco de planificación
y prácticas XP para el diseño, la implementación y las pruebas. Se reutilizaron
las historias de usuario y los casos de prueba de la Unidad~1, de modo que la
versión en JavaScript se valida con los mismos criterios que la versión en Go.

\subsection{Roles}
\begin{table}[H]
\centering
\caption{Distribución de roles y responsabilidades}
\begin{tabular}{ll}
\toprule
\textbf{Rol} & \textbf{Responsable}\\
\midrule
Product Owner & Ing. César Villacís Silva, MSc.\\
Scrum Master & Diana Guerra\\
Desarrollo -- UnbeatableAI e interfaz web & Germán Cáceres\\
Desarrollo -- MediumAI & Paul Moreno\\
Desarrollo -- UnbeatableAI $\rightarrow$ MinMax & Diana Guerra\\
\bottomrule
\end{tabular}
\end{table}

\subsection{Product Backlog}
\begin{longtable}{p{1.1cm}p{8.6cm}p{1.2cm}p{1.6cm}}
\caption{Historias de usuario}\\
\toprule
\textbf{HU} & \textbf{Historia} & \textbf{Pts} & \textbf{Sprint}\\
\midrule\endfirsthead
\toprule
\textbf{HU} & \textbf{Historia} & \textbf{Pts} & \textbf{Sprint}\\
\midrule\endhead
HU01 & Como usuario, quiero ver un tablero 3$\times$3 en el navegador, para jugar una partida de Tres en Raya. & 3 & 1\\
HU02 & Como usuario, quiero marcar una casilla y recibir la respuesta de la JCC, para alternar los turnos de la partida. & 5 & 1\\
HU03 & Como usuario, quiero enfrentar a UnbeatableAI con reglas de victoria y bloqueo, para jugar contra una estrategia que aproveche oportunidades y responda a mis amenazas. & 8 & 1\\
HU04 & Como usuario, quiero conocer el ganador o el empate, para identificar el resultado de la partida. & 3 & 1\\
HU05 & Como usuario, quiero elegir la modalidad de IA, para jugar contra la estrategia seleccionada. & 3 & 2\\
HU06 & Como usuario, quiero animaciones y sonidos durante la partida, para recibir retroalimentación visual y auditiva. & 5 & 2\\
HU07 & Como estudiante, quiero ver los pesos y sumas de las IA numéricas, para comprender sus decisiones. & 3 & 2\\
HU08 & Como usuario, quiero consultar un marcador por modalidad, para conocer mis partidas ganadas, empatadas y perdidas. & 2 & 2\\
HU09 & Como usuario, quiero jugar contra MediumAI, para enfrentar su estrategia de centro, victoria, bloqueo y selección aleatoria. & 5 & 3\\
HU10 & Como usuario, quiero jugar contra Minimax con las decisiones de la variante C\#, para enfrentar una búsqueda completa sin derrotas. & 8 & 3\\
\bottomrule
\end{longtable}

\subsection{Sprints}
\begin{description}
  \item[Sprint 1 -- Núcleo del juego.] Migración del método numérico a JavaScript
        (\texttt{unbeatable-ai.js}), tablero (\texttt{board.js}), despacho de clics
        (\texttt{form.js}) y pruebas de los Casos 1--3.
  \item[Sprint 2 -- Experiencia web.] Interfaz con HTML, CSS y DOM: tablero con
        trazos animados, línea ganadora, sonidos sintetizados con Web Audio,
        inspector de IA, marcador, bitácora y tarjeta de resultado.
  \item[Sprint 3 -- MediumAI y Minimax.] Migración de MediumAI y de Minimax con el
        flujo de su formulario, pruebas automatizadas, comparación con el C\# y
        validación en el navegador.
\end{description}
Se completó HU10 con 2\,097 decisiones coincidentes con el C\# y 569 finales sin derrotas.

\subsection{Prácticas XP aplicadas}
\begin{itemize}
 \item \textbf{Diseño simple:} un módulo por modalidad, sin frameworks, con los
       mismos métodos en los tres controladores.
 \item \textbf{Pruebas automatizadas:} 24 pruebas de primer nivel con el ejecutor
       nativo de Node.js (\texttt{node:test}), sin bibliotecas externas.
 \item \textbf{Integración incremental:} la interfaz se conectó a \texttt{Form} y
       \texttt{Board}, manteniendo la estrategia de IA separada del DOM.
 \item \textbf{Reutilización de la Unidad~1:} los mismos identificadores de prueba
       (U, D y M) y la misma referencia SHA-256 permiten comparar ambas versiones.
\end{itemize}

\subsection{Diseño}
La aplicación se organiza en dos capas dentro del navegador:
\begin{itemize}
  \item \textbf{Lógica (\texttt{src/game}):} los controladores trabajan sobre
        \texttt{Board}, que reemplaza a \texttt{ListButtons}, \texttt{GroupBox} y
        \texttt{MessageBox}, de modo que el código sigue casi línea por línea al C\#.
        No accede al DOM, por eso las pruebas lo ejecutan en Node.js.
  \item \textbf{Interfaz (\texttt{src/ui}):} \texttt{app.js} dibuja el tablero y
        reproduce con pausas los eventos que deja cada jugada; no decide ninguna jugada.
\end{itemize}

\begin{table}[H]
\centering\small
\caption{Estructura del proyecto}
\begin{tabular}{p{4.6cm}p{9cm}}
\toprule
\textbf{Archivo} & \textbf{Responsabilidad}\\
\midrule
\texttt{index.html}, \texttt{css/styles.css} & Estructura y estilo de la página, adaptable a celular\\
\texttt{src/game/constants.js} & \texttt{Player}, \texttt{GameState}, índices, líneas y catálogo de modos\\
\texttt{src/game/board.js} & Tablero visible y registro de eventos\\
\texttt{src/game/form.js} & \texttt{frmTicTacToe}: nivel, despacho de clics y reinicio\\
\texttt{src/game/*-ai.js} & Un controlador por modalidad\\
\texttt{src/ui/app.js}, \texttt{sound.js} & Dibujo, animación, marcador, bitácora y sonido\\
\texttt{server.js} & Servidor estático sin dependencias\\
\texttt{test/*.test.js} & Pruebas U01--U08, D01--D08 y M01--M08\\
\bottomrule
\end{tabular}
\end{table}

\subsection{Desarrollo}
La migración de C\# a JavaScript estableció las siguientes correspondencias:

\begin{table}[H]
\centering
\caption{Correspondencia C\# $\rightarrow$ JavaScript}
\small
\begin{tabular}{ll}
\toprule
\textbf{C\# (original)} & \textbf{JavaScript}\\
\midrule
\texttt{ListButtons[i].Text} & \texttt{board.buttons[i]}\\
\texttt{grbTicTacToe.Enabled} & \texttt{board.enabled}\\
\texttt{MessageBox.Show(...)} & \texttt{board.showMessage(...)} $\rightarrow$ evento\\
\texttt{enum Player} / \texttt{enum GameState} & \texttt{Object.freeze(\{...\})}\\
\texttt{mBoardState} (\texttt{Player[]}) & \texttt{boardState}: arreglo de 9\\
\texttt{mArrayMoves} (\texttt{int[]}) & \texttt{arrayMoves}: arreglo de 9\\
\texttt{CENTER\_INDEX}, \texttt{CORNER\_/SIDE\_INDICES} & \texttt{CENTER\_INDEX}, \texttt{CORNER\_/SIDE\_INDICES}\\
\texttt{new Random().Next(0, 9)} & \texttt{Math.floor(random() * 9)}\\
\texttt{public string Flag \{ get; \}} & \texttt{get flag()}\\
\texttt{static bool ControlPlayer} & \texttt{form.controlPlayer}\\
\texttt{btnTicTacToeN\_Click} & \texttt{form.click(index)}\\
\bottomrule
\end{tabular}
\end{table}

El generador aleatorio se recibe en el constructor (\texttt{Math.random} por
defecto), lo que permite inyectar uno controlado sin cambiar la lógica. En la IA
fuerte se incorporó una comprobación de disponibilidad en los \texttt{do/while}
aleatorios: si no queda ninguna casilla del grupo pedido, el original entra en un
bucle infinito, que en el navegador congelaría la pestaña (RNF07).

\subsection{Pruebas}
Se realizaron pruebas funcionales y de integración el 19 de noviembre de 2026.
Desde la carpeta \texttt{UNIDAD II/Projectos/TickTackToe-JavaScript} se ejecutaron:
\begin{lstlisting}
npm test                  # node --test
npm run test:coverage     # node --test --experimental-test-coverage
\end{lstlisting}
Se aprobaron 24 pruebas de primer nivel (8 por modalidad) con 77 subcasos, en
total 101 pruebas sin fallos. Las tablas siguientes presentan los resultados de
una ejecución por archivo; las exploraciones aleatorias (D08 y U08) varían entre
ejecuciones.

\subsubsection{Pruebas automatizadas de MediumAI}
\begin{longtable}{p{1cm}p{3.4cm}p{3.3cm}p{3.4cm}p{1.6cm}}
\caption{Objetivos, escenarios y resultados de MediumAI}\\
\toprule ID & Objetivo y procedimiento & Esperado & Obtenido & Estado\\\midrule\endfirsthead
\toprule ID & Objetivo y procedimiento & Esperado & Obtenido & Estado\\\midrule\endhead
D01 & Inicializar y reiniciar controlador y tablero. & Pesos vacíos, turno X y tablero habilitado. & Estado y tablero restablecidos. & PASS\\
D02 & X marca 0 con el centro libre. & O toma el centro (4). & O en 4; turno X. & PASS\\
D03 & Intentar jugar en el centro ocupado. & Sin cambios en tablero, pesos ni eventos. & Jugada ocupada ignorada. & PASS\\
D04 & Probar las ocho líneas, suma 6/2 y cada casilla vacía. & Ganar o bloquear en el espacio correcto. & 48 subcasos aprobados. & PASS\\
D05 & Presentar una victoria y un bloqueo posibles. & Priorizar suma 6. & Selecciona la victoria en 8. & PASS\\
D06 & Sin suma 6/2 y una casilla libre. & Elegir al azar entre libres. & Elige la casilla libre 4. & PASS\\
D07 & Comprobar victoria de X y empate. & Mensaje y tablero deshabilitado. & Ambos finales correctos. & PASS\\
D08 & Explorar 30 árboles con respuestas aleatorias. & Registrar finales y derrotas de la IA. & 15\,934 finales; 606 derrotas (3,80\%). & PASS\\
\bottomrule
\end{longtable}

\begin{figure}[H]
\lstinputlisting[style=consola]{Img/consola-medium-ai.txt}
\caption{Salida de consola de las pruebas de MediumAI.}
\label{fig:consola-medium}
\end{figure}

\subsubsection{Pruebas automatizadas de UnbeatableAI}
\begin{longtable}{p{1cm}p{3.4cm}p{3.3cm}p{3.4cm}p{1.6cm}}
\caption{Objetivos, escenarios y resultados de UnbeatableAI}\\
\toprule ID & Objetivo y procedimiento & Esperado & Obtenido & Estado\\\midrule\endfirsthead
\toprule ID & Objetivo y procedimiento & Esperado & Obtenido & Estado\\\midrule\endhead
U01 & Primera respuesta a X en 0. & O ocupa el centro. & O en 4. & PASS\\
U02 & Humano ocupa el centro; 50 repeticiones. & O elige una esquina. & Las 50 respuestas son esquinas. & PASS\\
U03 & Caso 1: X4, O0; humano marca 2. & Bloqueo en 6. & O en 6. & PASS\\
U04 & Caso 2: O en 2 y 8; humano marca 6. & O gana en 5 y deshabilita el tablero. & Jugada, mensaje y final correctos. & PASS\\
U05 & Amenaza y victoria simultáneas; X marca 5. & Priorizar la victoria diagonal. & O en 8. & PASS\\
U06 & X en lados 1 y 3; luego 5. & Primera regla de esquina: 0. & O en 0. & PASS\\
U07 & X0, O4; X8; 200 repeticiones. & O en 1, 5 o 7; nunca 3. & Restricción original conservada. & PASS\\
U08 & Recorrer las opciones humanas en 30 árboles con IA aleatoria. & Partidas terminan; derrotas de la IA no mayores al 3\%. & 15\,006 finales, 151 derrotas (1,01\%); 0 bloqueos. & PASS\\
\bottomrule
\end{longtable}

\begin{figure}[H]
\lstinputlisting[style=consola]{Img/consola-unbeatable-ai.txt}
\caption{Salida de consola de las pruebas de UnbeatableAI.}
\label{fig:consola-unbeatable}
\end{figure}

\subsubsection{Pruebas automatizadas de Minimax}
\begin{longtable}{p{1cm}p{3.4cm}p{3.3cm}p{3.4cm}p{1.6cm}}
\caption{Objetivos, escenarios y resultados de Minimax}\\
\toprule ID & Objetivo y procedimiento & Esperado & Obtenido & Estado\\\midrule\endfirsthead
\toprule ID & Objetivo y procedimiento & Esperado & Obtenido & Estado\\\midrule\endhead
M01 & Reiniciar tras victoria de la IA; comparar con el constructor. & Matriz vacía, turno X y tablero habilitado. & Datos y controles reiniciados. & PASS\\
M02 & X0; índices inválidos, doble turno y casillas ocupadas. & Rechazo de ilegales; O4 y pesos 1/3. & Turnos, pesos y eventos correctos. & PASS\\
M03 & Tablero vacío, aperturas, victoria, bloqueo y tablero lleno. & Índices 0, 4, 0, 2, 2, $-1$; restauración íntegra. & Índices y matriz correctos. & PASS\\
M04 & Ocho líneas con X y O; profundidades 0, 1 y 4. & $+10/-10$ y ajuste por profundidad. & Líneas, puntuaciones y orden correctos. & PASS\\
M05 & Completar finales e intentar más movimientos. & Mensajes originales y ninguna jugada posterior. & Cuatro escenarios terminales correctos. & PASS\\
M06 & Ramificar todos los movimientos humanos contra la IA. & Ninguna victoria humana. & 569 finales: 386 victorias de la IA, 183 empates, 0 derrotas. & PASS\\
M07 & Enumerar estados legales; comparar decisiones con el C\#. & Mismos conteos y SHA-256. & 5\,478 estados; 2\,097 decisiones coincidentes. & PASS\\
M08 & \texttt{Form} con nivel 3: click, reinicio y cambio de nivel. & Despacho, rechazo y reinicio correctos. & Sin O extra tras final humano; reinicio completo. & PASS\\
\bottomrule
\end{longtable}

M07 enumera los estados legales empezando con X y corta cada rama al ganar o
llenar el tablero. Deduplica 5\,478 tableros y compara las 2\,097 posiciones no
terminales con turno de O. Ordena los tableros y concatena nueve caracteres,
barra vertical, índice devuelto y salto de línea; calcula SHA-256 en UTF-8. La
referencia se obtuvo en la Unidad~1 al compilar y ejecutar el C\# original. La
versión JavaScript produjo el mismo resumen, lo que demuestra que toma las mismas
decisiones que el C\# y que la versión Go en todo ese dominio:
\begin{lstlisting}
fb6a9200c81c724a24f04ddad024e88659df8a3837e3c2c4c5fb9d6c84d7b31f
\end{lstlisting}

\begin{figure}[H]
\lstinputlisting[style=consola]{Img/consola-minimax-ai.txt}
\caption{Salida de consola de las pruebas de Minimax, con la comparación SHA-256.}
\label{fig:consola-minimax}
\end{figure}

\subsubsection{Integración en el navegador}
Para verificar la integración se inició el servidor con \texttt{npm start} y se
jugaron partidas reales en Google Chrome 157, controlado mediante un script que
pulsa las teclas de las casillas en la interfaz.
\begin{longtable}{p{1cm}p{3.4cm}p{3.3cm}p{3.4cm}p{1.6cm}}
\caption{Integración en el navegador}\\
\toprule ID & Objetivo y procedimiento & Esperado & Obtenido & Estado\\\midrule\endfirsthead
\toprule ID & Objetivo y procedimiento & Esperado & Obtenido & Estado\\\midrule\endhead
I01 & Solicitar la página, un módulo, la hoja de estilos y una ruta fuera de la carpeta. & 200 con el tipo correcto; 404 fuera de la carpeta. & 200 en los tres archivos; 404 al salir de la carpeta con \texttt{/../}. & PASS\\
I02 & Seleccionar cada modo y jugar dos turnos. & Inspector en los modos numéricos y oculto en Minimax. & Correcto en los tres modos. & PASS\\
I03 & Jugar hasta el final en cada modo. & Tarjeta con el mensaje original de la variante. & «¡Ganó X!», «¡Ganó O!», «AI Wins!» e «It's a Draw!». & PASS\\
I04 & Revisar el tablero tras cada final. & Ninguna casilla habilitada. & 0 casillas habilitadas en los tres modos. & PASS\\
I05 & Abrir la página a 390 px de ancho. & Sin desplazamiento horizontal. & 0 px de desborde. & PASS\\
I06 & Registrar errores de la página durante todas las partidas. & Ningún error. & Ningún error en consola. & PASS\\
\bottomrule
\end{longtable}

\subsubsection{Validación visual en el navegador}
Las figuras~\ref{fig:captura-modos}--\ref{fig:captura-resultados} presentan la
pantalla principal, una partida de cada modalidad con su tablero final y la
tarjeta de resultado de cada variante.

\begin{figure}[H]
\centering
\begin{subfigure}[b]{0.72\linewidth}\centering
\captura[height=5.4cm]{js-captura-modos.png}
\caption{Escritorio.}
\end{subfigure}\hfill
\begin{subfigure}[b]{0.26\linewidth}\centering
\captura[height=5.4cm,trim=0 488 0 0,clip]{js-captura-celular.png}
\caption{Celular (390 px).}
\end{subfigure}
\caption{Pantalla principal con MediumAI, UnbeatableAI y Minimax disponibles.}
\label{fig:captura-modos}
\end{figure}

\begin{figure}[H]
\centering
\begin{subfigure}[t]{0.49\linewidth}\centering
\captura[width=0.99\linewidth,trim=170 0 170 0,clip]{js-captura-medium.png}
\caption{Partida en curso.}
\end{subfigure}\hfill
\begin{subfigure}[t]{0.49\linewidth}\centering
\captura[width=0.99\linewidth,trim=170 0 170 0,clip]{js-final-medium.png}
\caption{Victoria de X: línea ganadora y tablero bloqueado.}
\end{subfigure}
\caption{Partida de MediumAI, con pesos y sumas en el inspector.}
\label{fig:captura-medium}
\end{figure}

\begin{figure}[H]
\centering
\begin{subfigure}[t]{0.49\linewidth}\centering
\captura[width=0.99\linewidth,trim=170 0 170 0,clip]{js-captura-unbeatable.png}
\caption{Partida en curso, con suma 6 en la columna 2.}
\end{subfigure}\hfill
\begin{subfigure}[t]{0.49\linewidth}\centering
\captura[width=0.99\linewidth,trim=170 0 170 0,clip]{js-final-unbeatable.png}
\caption{Victoria de O: línea ganadora y tablero bloqueado.}
\end{subfigure}
\caption{Partida de UnbeatableAI, con la regla numérica observable en el inspector.}
\label{fig:captura-unbeatable}
\end{figure}

\begin{figure}[H]
\centering
\begin{subfigure}[t]{0.49\linewidth}\centering
\captura[width=0.99\linewidth,trim=170 0 170 0,clip]{js-captura-minimax.png}
\caption{Partida en curso.}
\end{subfigure}\hfill
\begin{subfigure}[t]{0.49\linewidth}\centering
\captura[width=0.99\linewidth,trim=170 0 170 0,clip]{js-final-minimax.png}
\caption{Victoria de la IA: línea ganadora y tablero bloqueado.}
\end{subfigure}
\caption{Partida de Minimax, con tablero y bitácora y sin inspector numérico.}
\label{fig:captura-minimax}
\end{figure}

\begin{figure}[H]
\centering
\begin{subfigure}[t]{0.49\linewidth}\centering
\captura[width=0.99\linewidth,trim=560 280 560 280,clip]{js-resultado-medium.png}
\caption{MediumAI: «¡Ganó X!».}
\end{subfigure}\hfill
\begin{subfigure}[t]{0.49\linewidth}\centering
\captura[width=0.99\linewidth,trim=560 280 560 280,clip]{js-resultado-unbeatable.png}
\caption{UnbeatableAI: «¡Ganó O!».}
\end{subfigure}\\[6pt]
\begin{subfigure}[t]{0.49\linewidth}\centering
\captura[width=0.99\linewidth,trim=560 280 560 280,clip]{js-resultado-minimax-victoria.png}
\caption{Minimax: «AI Wins!».}
\end{subfigure}\hfill
\begin{subfigure}[t]{0.49\linewidth}\centering
\captura[width=0.99\linewidth,trim=560 280 560 280,clip]{js-resultado-minimax.png}
\caption{Minimax: «It's a Draw!».}
\end{subfigure}
\caption{Tarjeta de resultado de cada modalidad, con el mensaje y el título originales de su variante C\#.}
\label{fig:captura-resultados}
\end{figure}

Las pruebas en el navegador documentan partidas concretas y no sustituyen una
suite automatizada de extremo a extremo.

% =====================================================================
\section{Análisis de Resultados: Métricas de Software}
% =====================================================================

\subsection{Líneas de código}
Se contabilizaron las líneas de código el 19 de noviembre de 2026 con un script
de Node.js y el mismo criterio de la Unidad~1: una línea sin caracteres no blancos
es vacía; si, al retirar comentarios respetando cadenas, queda código, es LOC; si
solo queda comentario, es comentario. Una línea con código y comentario cuenta
una vez como LOC. En JavaScript se cuentan como funciones las declaraciones, los
métodos, los constructores y los \emph{getters}.

\begin{table}[H]
\centering\small
\caption{Tamaño por alcance de código fuente}
\begin{tabular}{lrrrrr}
\toprule Alcance & Archivos & LOC & Coment. & Blanco & Total\\\midrule
Lógica (\texttt{src/game}) & 6 & 554 & 82 & 74 & 710\\
Interfaz (\texttt{src/ui}, HTML y CSS) & 4 & 672 & 19 & 61 & 752\\
Servidor estático & 1 & 30 & 2 & 2 & 34\\
Pruebas (\texttt{test}) & 4 & 422 & 11 & 34 & 467\\
\textbf{Total} & \textbf{15} & \textbf{1\,678} & \textbf{114} & \textbf{171} & \textbf{1\,963}\\
\bottomrule
\end{tabular}
\end{table}

\begin{table}[H]
\centering\small
\caption{Controladores por modalidad: C\#, Go (Unidad 1) y JavaScript}
\begin{tabular}{lrrrrr}
\toprule Archivo & LOC & Coment. & Blanco & Total & Func.\\\midrule
C\# MediumAIController & 315 & 83 & 39 & 437 & 13\\
Go medium\_ai.go & 240 & 4 & 14 & 258 & 14\\
JS medium-ai.js & 102 & 12 & 16 & 130 & 16\\\midrule
C\# UnbeatableAIController (fuerte) & 417 & 83 & 39 & 539 & 13\\
Go unbeatable\_ai.go & 223 & 39 & 26 & 288 & 17\\
JS unbeatable-ai.js & 170 & 26 & 21 & 217 & 17\\\midrule
C\# UnbeatableAIController (Minimax) & 120 & 3 & 15 & 138 & 4\\
Go minimax\_ai.go & 228 & 17 & 17 & 262 & 15\\
JS minimax-ai.js & 149 & 16 & 17 & 182 & 15\\
\bottomrule
\end{tabular}
\end{table}
La reducción en MediumAI y UnbeatableAI se debe principalmente a que las 16 ramas
\texttt{if/else if} de las sumas 6 y 2 se expresan como dos recorridos de
\texttt{WINNING\_LINES}. En Minimax, el archivo JavaScript reúne además el flujo
migrado de \texttt{frmTicTacToe}, por eso supera al controlador C\#. El tamaño
del código se interpreta como una métrica descriptiva, no de rendimiento.

\subsection{Otras métricas}
\begin{table}[H]
\centering\small
\caption{Métricas funcionales por modalidad}
\begin{tabular}{p{4.6cm}rrr}
\toprule Métrica & MediumAI & UnbeatableAI & Minimax\\\midrule
Archivos de controlador / prueba & 1 / 1 & 1 / 1 & 1 / 1\\
Funciones del controlador & 16 & 17 & 15\\
Pruebas de primer nivel aprobadas & 8 / 8 & 8 / 8 & 8 / 8\\
Finales explorados por sus pruebas & 15\,934 & 15\,006 & 569\\
Derrotas de la IA & 606 & 151 & 0\\
Porcentaje de derrotas de la IA & 3,80\% & 1,01\% & 0\%\\
Cobertura de líneas & 97,69\% & 93,55\% & 100\%\\
Dependencias externas & 0 & 0 & 0\\
\bottomrule
\end{tabular}
\end{table}

Se obtuvo una cobertura de \textbf{97,61\% de líneas}, 96,90\% de ramas y 95,52\%
de funciones en \texttt{src/game}. Las líneas sin cubrir de UnbeatableAI coinciden
con los hallazgos de la sección~\ref{sec:hallazgos}: la protección contra el
bloqueo nunca se activa y la regla del centro a partir de la segunda jugada nunca
se ejecuta. El resto son \texttt{flag} y \texttt{weights}, que solo usa la
interfaz, y ramas defensivas equivalentes a las del original.

\begin{figure}[H]
\lstinputlisting[style=consola]{Img/consola-cobertura.txt}
\caption{Cobertura de la lógica del juego con el ejecutor nativo de Node.js.}
\label{fig:consola-cobertura}
\end{figure}

\subsubsection{Tiempo de decisión}
Se midió \texttt{form.click(0)} en Node.js~24.19.0, sobre Windows~11
(10.0.26300) con un procesador AMD Ryzen~5 7520U. La medición incluye la jugada
del usuario y la respuesta de la IA, y excluye las pausas que la interfaz añade
a propósito para animar. Se descartaron 5 ejecuciones de calentamiento por modo
y se tomaron 30 muestras: mediana = promedio de las posiciones 15 y 16;
P95 = posición 29.
\begin{table}[H]
\centering\small
\caption{Tiempo de la primera respuesta de la IA, en milisegundos}
\begin{tabular}{lrrrrr}
\toprule Modo & Muestras & Mínimo & Mediana & P95 & Máximo\\\midrule
MediumAI & 30 & 0,004 & 0,005 & 0,012 & 0,022\\
UnbeatableAI & 30 & 0,004 & 0,006 & 0,017 & 0,441\\
Minimax & 30 & 2,744 & 2,776 & 3,332 & 3,481\\
\bottomrule
\end{tabular}
\end{table}
Todas las muestras están muy por debajo de los 50~ms de RNF02. Minimax es el más
lento porque explora todas las continuaciones legales, pero incluso en el tablero
casi vacío responde en menos de 4~ms. La interfaz espera 550~ms antes de mostrar
la jugada de la JCC para que el usuario la perciba.

\subsection{Comportamiento de la variante UnbeatableAI}\label{sec:hallazgos}
La migración reprodujo los mismos comportamientos de la estrategia numérica que
se identificaron en la Unidad~1:

\begin{enumerate}
  \item \textbf{Derrotas observadas.} Se registraron \textbf{151 derrotas de la JCC
        en 15\,006 finales (1,01\%)} durante el recorrido con muestreo. Cuando
        ninguna regla aplica, la JCC juega al azar y a veces deja al usuario un
        \emph{tenedor} (dos amenazas a la vez). Ejemplo: X en 0, O en 4, X en 7,
        O en 2 (al azar), X en 6; la JCC solo puede tapar una de las amenazas
        (0--3--6 y 6--7--8).
  \item \textbf{La regla de suma 5 excluye el lado 3.} La condición original del
        lado aleatorio compara \texttt{SIDE\_INDICES[0]} dos veces, así que el
        lado 3 nunca se elige en esa regla. La versión JavaScript conserva el
        error y la prueba U07 lo verifica.
  \item \textbf{Posible bloqueo.} Los \texttt{do/while} aleatorios no tienen salida
        si no queda casilla válida del grupo. En los 15\,006 finales se registraron
        \textbf{0 casos}; aun así se mantiene la comprobación de disponibilidad.
  \item \textbf{Regla del centro sin efecto.} A partir de la segunda jugada,
        \texttt{playNPC} revisa si el centro está libre; nunca lo está a esa altura.
        La cobertura lo confirma: esas líneas no se ejecutaron en ninguna prueba.
\end{enumerate}

Con \textbf{Minimax} se obtuvieron cero derrotas en los 569 finales explorados
frente a todos los movimientos humanos legales.

% =====================================================================
\section{Conclusiones y Recomendaciones}
% =====================================================================
\subsection{Conclusiones}
\begin{itemize}
 \item Se migraron las tres modalidades del Tres en Raya a JavaScript para la web,
       sin frameworks ni dependencias. La separación entre la lógica
       (\texttt{src/game}) y la interfaz (\texttt{src/ui}) permitió probar las
       reglas en Node.js y jugar en el navegador con el mismo código.
 \item Minimax en JavaScript toma las mismas 2\,097 decisiones que el C\# original,
       con el mismo SHA-256 obtenido en la Unidad~1; en 569 finales no perdió ninguna partida.
 \item MediumAI y UnbeatableAI conservan sus reglas, incluidos los comportamientos
       del original: UnbeatableAI perdió el 1,01\% de los finales explorados por
       los tenedores que dejan sus jugadas al azar.
 \item Se aprobaron 24 pruebas de primer nivel (101 con subcasos), con 97,61\% de
       cobertura de líneas. En el navegador se jugaron partidas de los tres modos,
       sin errores y con el tablero bloqueado al terminar.
 \item Frente a la versión Go y Angular, esta versión elimina el servidor de
       lógica y la API: la respuesta de la IA es inmediata (menos de 4~ms) y
       el proyecto se ejecuta solo con Node.js y un navegador.
\end{itemize}

\subsection{Recomendaciones}
\begin{itemize}
 \item Incorporar una suite de pruebas de navegador de extremo a extremo.
 \item Mantener la referencia SHA-256 y las pruebas M06 y M07 para detectar
       cambios de comportamiento en futuras modificaciones.
 \item Si se desea una IA verdaderamente imbatible con el método numérico, evaluar
       como mejora aparte reemplazar las jugadas al azar, ya que hacerlo cambiaría
       la lógica del original.
\end{itemize}

% =====================================================================
\begin{thebibliography}{9}
\bibitem{guia-lab} C. Villacís Silva, \emph{Informe de laboratorio Desarrollo de Videojuegos}, indicaciones del Proyecto 3.
\bibitem{minimax-docente} Material del docente, \emph{Tic-Tac-Toe-Version-MiniMax}: \texttt{UnbeatableAIController.cs}, \texttt{frmTicTacToe.cs} y \texttt{minmax-ia-explicacion.docx}, incluidos en \texttt{UNIDAD I/Material}.
\bibitem{ieee830} IEEE, \emph{IEEE Recommended Practice for Software Requirements Specifications}, IEEE Std 830-1998.
\bibitem{villacis} C. Villacís, W. Fuertes, A. Bustamante et al., \emph{Optimización del Juego del Tres en Raya con Niveles de Dificultad, Utilizando Técnicas Heurísticas de Inteligencia Artificial}, Universidad de las Fuerzas Armadas -- ESPE.
\bibitem{mdn} Mozilla, \emph{MDN Web Docs: JavaScript}, \url{https://developer.mozilla.org/docs/Web/JavaScript}.
\bibitem{node} OpenJS Foundation, \emph{Node.js Documentation: Test runner}, \url{https://nodejs.org/api/test.html}.
\bibitem{scrum} K. Schwaber y J. Sutherland, \emph{La Guía de Scrum}, 2020.
\bibitem{xp} K. Beck, \emph{Extreme Programming Explained: Embrace Change}, Addison-Wesley, 2004.
\end{thebibliography}

\end{document}
